\chapter{The Command-Line Toolbox}
\label{ch:toolbox}

\standfirst{Nine things \ct{st80} can tell you about an image or a source tree
without opening a window. Most of them exist because a question came up during
development and grepping was not good enough.}

\section{\texttt{-syntax}: does it compile?}

\begin{sh}
./st80 -syntax lib/MyStuff/*.st
./st80 -syntax -profile profiles/mystuff.profile
\end{sh}

Compiles every method in the named files and reports failures, without
building an image or running anything. The fastest check there is, and the
right thing to put in an editor's save hook.

It is a \emph{grammar} check and deliberately nothing more: every identifier
resolves, and its stand-in literals are shared rather than allocated, so it
sees neither an undefined global nor a method that overruns the 63-literal
frame. Those are the bootstrap's business, and \ct{-tests} is the next step
in the workflow below for exactly that reason.

\section{\texttt{-primitives}: what does this source ask of the VM?}

\begin{sh}
./st80 -primitives pharo/Collections-Unordered/*.st
\end{sh}

Lists every primitive number the source declares. This turns "will this port
run?" from an unbounded question into a finite checklist: the primitives it
names either exist here or they do not.

\section{\texttt{-census}: what is in this image?}

\begin{sh}
./st80 -census st80.image
\end{sh}

\begin{plain}
live objects     : 15607
pointer objects  : 5267
non-pointer      : 10340
words in objects : 1028771
sum of refcounts : 49585

reference count histogram (count: objects)
 10632 1
   800 2
   ...
\end{plain}

A summary of the object memory. Useful for noticing that something is
retaining far more than it should.

\section{\texttt{-classes} and \texttt{-methods}}

\begin{sh}
./st80 -classes st80.image | grep -i collection
./st80 -methods st80.image | grep 'inject:into:'
\end{sh}

Every class and every method, in the format of the Xerox reference dumps.
These exist because those dumps are an oracle: \ct{class.oops} lists 446
classes and \ct{method.oops} lists 4,494 methods, and this system reproduces
both exactly. As a day-to-day tool it is grep over the whole system.

\section{\texttt{-disasm}: what did the compiler emit?}

\begin{sh}
./st80 -disasm st80.image OrderedCollection add:
\end{sh}

\begin{plain}
OrderedCollection>>add:
  header      : ... (flag 0, 1 temporaries, 3 literals)
  literal 0   : ...
  bytecodes   :
      ...: 112          Push Receiver
      ...:  16          Push Temporary 0
      ...
\end{plain}

For when you want to know whether \ct{ifTrue:} really was inlined, or whether
a send is going through the special-selector fast path.

\section{\texttt{-inspect}: one object, by pointer}

\begin{sh}
./st80 -inspect st80.image 1a4c
\end{sh}

Describes one object, given its oop in hex. Mostly useful with \ct{-census}
and \ct{-trace2} output, which name objects that way.

\section{\texttt{-trace2} and \texttt{-trace3}: the oracle}

\begin{sh}
make OM=bb
./st80 -trace2 oracle/VirtualImage 611
\end{sh}

The Xerox execution traces, byte for byte. These need the 1983 image, which
carries no licence grant from anyone and so is not distributed---the suite
skips without it and everything else still passes.

If you change the interpreter, this is the test that matters.

\section{\texttt{-inject} and \texttt{-screenshot}: driving the interface}

\begin{sh}
./st80 -inject "m 320 240; w 20; d 129; w 40; u 129; w 60" \
       -screenshot /tmp/shot.pbm -run st80.image
\end{sh}

\begin{center}
\small
\begin{tabular}{@{}ll@{}}
\toprule
\ctp{m X Y} & move the pointer\\
\ctp{d CODE} \ctp{u CODE} & press and release (128 right, 129 middle, 130 left)\\
\ctp{k CODE} & type a key\\
\ctp{w N} & wait N bytecode slices\\
\ctp{x} & post a window-exposed event\\
\bottomrule
\end{tabular}
\end{center}

The injected events go through the same queue SDL's own do, which is the
point: a test that drove a private queue would prove nothing about the one the
image reads.

\section{\texttt{-eval}, \texttt{-tests}, \texttt{-doctests}}

Covered in Chapters~\ref{ch:building} and~\ref{ch:tests}. The three you will
use daily.

\section{A workflow}

\begin{sh}
# fast: does it compile?
./st80 -syntax lib/MyStuff/*.st

# does it work?
./st80 -bootstrap -profile profiles/mystuff.profile -tests

# do the documented examples still hold?
./st80 -bootstrap -profile profiles/mystuff.profile \
       -doctests lib/MyStuff/Account.class.st

# does the whole system still work?
make test

# and before anything touching threads
make clean && make TSAN=1 test
\end{sh}

The first takes under a second, the second about ten, and \ct{make test} a few
minutes. Run them in that order and you will rarely wait for the slow one to
tell you something the fast one could have.
